% Generated draft: requires mathematical and source review.
\documentclass[11pt]{article}
\usepackage[margin=25mm]{geometry}
\usepackage{fontspec}
\setmainfont{Latin Modern Roman}
\usepackage{amsmath,amssymb,mathtools,unicode-math}
\setmathfont{Latin Modern Math}
\usepackage{enumitem,needspace,xurl,hyperref}
\setlength{\parindent}{0pt}
\setlength{\parskip}{5pt}
\setlength{\emergencystretch}{3em}
\newcommand{\N}{\mathbb{N}}
\newcommand{\Z}{\mathbb{Z}}
\newcommand{\Q}{\mathbb{Q}}
\newcommand{\R}{\mathbb{R}}
\newcommand{\C}{\mathbb{C}}
\newcommand{\F}{\mathbb{F}}
\newcommand{\A}{\mathbb{A}}
\newcommand{\PP}{\mathbb{P}}
\newcommand{\E}{\mathbb{E}}
\newcommand{\eps}{\varepsilon}
\newcommand{\dd}{\,\mathrm{d}}
\begin{document}
\section*{UnsolvedMath ID: 30004994}
\textbf{Source problem code:} catalogue entry
\subsection{Definitions and mathematical statement}
The supplied record asks about an algebraic action on a compact metrizable space. Under the conventional interpretation, this means that \(X\) is a compact metrizable topological group and that each element of \(\Gamma\) acts by a continuous group automorphism. The record does not explicitly specify whether \(X\) must be abelian. The formulation below uses this interpretation of algebraic action; its intended scope requires review.

Let \(\Gamma\) be a countably infinite discrete amenable group, and let \(\Gamma\curvearrowright X\) be such an algebraic action. Fix a metric \(d\) compatible with the topology of \(X\). The action is expansive if there exists \(c>0\) such that
\[
\forall x,y\in X,\qquad \bigl[d(gx,gy)\leq c\text{ for every }g\in\Gamma\bigr]\Longrightarrow x=y.
\]
A nondiagonal asymptotic pair is a pair \((x,y)\in X^2\) with \(x\ne y\) such that
\[
\forall\varepsilon>0\ \exists K\subseteq\Gamma\text{ finite}\ \forall g\in\Gamma\setminus K,\qquad d(gx,gy)<\varepsilon.
\]
Thus asymptotic convergence here means convergence outside finite subsets of the whole group.

Interpret the unspecified entropy \(h\) in the record as topological entropy. Explicitly, choose a left F\o lner sequence \((F_n)_{n\geq1}\) of nonempty finite subsets of \(\Gamma\), so that
\[
\frac{|sF_n\mathbin{\triangle}F_n|}{|F_n|}\longrightarrow0\qquad(s\in\Gamma).
\]
For a finite open cover \(\mathcal U\) of \(X\), let \(N(\mathcal U)\) be its minimum subcover cardinality, and put
\[
h_{\mathrm{top}}(\Gamma\curvearrowright X)
=\sup_{\mathcal U}\lim_{n\to\infty}\frac{1}{|F_n|}
\log N\!\left(\bigvee_{g\in F_n}g^{-1}\mathcal U\right),
\]
where \(g^{-1}\mathcal U=\{\{x\in X:gx\in U\}:U\in\mathcal U\}\), the join consists of intersections of cover members, and the supremum ranges over finite open covers.

\textbf{Question under these conventions.} For every \(\Gamma\) and every expansive algebraic action as above, is
\[
h_{\mathrm{top}}(\Gamma\curvearrowright X)>0
\quad\Longleftrightarrow\quad
\exists x,y\in X\;\left[x\ne y\ \text{and}\ d(gx,gy)\longrightarrow0\text{ as }g\to\infty\right]?
\]
The interpretation of entropy, the algebraic structure required on \(X\), and the intended asymptotic convention are not explicitly defined in the supplied excerpt. This is therefore a qualified reconstruction requiring mathematical review.
\subsection{Short English statement}
For an expansive algebraic action of a countably infinite discrete amenable group, ask whether positive topological entropy is equivalent to the existence of two distinct points whose translates approach one another outside every sufficiently large finite subset of the group. The supplied record labels the problem partially solved; the precise algebraic scope and entropy convention require review.
\subsection{Sources}
\begin{itemize}
\item Local UnsolvedMath registry: supplied complete record for permanent ID \(30004994\), entitled \emph{Entropy and Nondiagonal Asymptotic Pairs}; exact source problem code \texttt{catalogue entry}. No individual proposer is supplied.
\item Supplied report citation: \emph{Groups and Dynamics: Topology, Measure, and Borel Structure} (2023), Oberwolfach Reports, pp.\ 107--169, DOI \texttt{10.4171/owr/2022/3}, \url{https://doi.org/10.4171/owr/2022/3}. The original extracted text includes a page marker 141. The registry separately gives proposed year 2022 and the tag \texttt{owr-year:2022}; the relationship between these dates and the citation year is not resolved here.
\item Chung--Li, \emph{Homoclinic groups, IE groups, and expansive algebraic actions}, \url{https://arxiv.org/abs/1103.1567}.
\item Meyerovitch, \emph{Pseudo-Orbit Tracing and Algebraic actions}, \url{https://arxiv.org/abs/1701.01318}.
\item Barbieri--Garc\'ia-Ramos--Li, \emph{Markovian properties of continuous group actions}, \url{https://arxiv.org/abs/1911.00785}.
\item The supplied literature assessment is dated exactly \texttt{2026-08-22}. This is supplied dated curation, not independently verified here. It labels the problem partially solved and reports that the equivalence is established for polycyclic-by-finite algebraic actions and related classes, while counterexamples and unresolved cases leave the general amenable expansive algebraic-action question unsettled. The record does not identify the precise relevant results or counterexample hypotheses. Its background expressly requests independent verification of time-sensitive status and cited claims.
\item The registry labels the statement \texttt{corrected\_verified} and reports that it was reconstructed and checked against the official Oberwolfach report. That is supplied verification metadata, not a verification performed for this draft.
\end{itemize}
\end{document}
